\chapter*{前言}

第一册正文讲清了 C++ 程序怎样经过编译、链接、装载、ABI 约定、指令执行和 Linux 工具链观察，最终成为一个正在运行的进程。但只读正文容易产生一种错觉：好像看懂概念就等于会做工程。真正的门槛在证据链里：你能不能拿一个自己编译出来的二进制，说明它是什么格式，入口在哪里，符号落在哪些 section，函数调用遵守什么 ABI，benchmark 结果是否被 checksum 保护，性能结论能否被反汇编、工具输出或测试复查。

本卷的贯穿项目叫 \term{Executable Evidence Kit}{可执行文件证据工具箱}。它不是一个玩具读文件程序，而是一条逐步扩展的研究项目：第一阶段读取字节并建立 checksum；第二阶段解析 ELF64 header；第三阶段解析 section table 和 string table；第四阶段解析 symbol table；第五阶段输出 report 和 CSV；第六阶段把 CLI、GTest/GMock、Google Benchmark、readelf/objdump/perf 观察和研究报告串起来。每一阶段都要回答同一个问题：这条结论来自哪里，怎样复查，当前工具还不能证明什么。

本卷补的是第一册最缺的两类训练：一类是章节习题，每章都要落到具体代码、命令和证据；另一类是研究项目，读者要从“照着跑”走到“能解释一个真实二进制”。因此本卷不会只列题目，也不会只贴模板。每个任务都先给问题背景，再说明最小实现、失败边界、测试要求和报告要求。

\begin{keyidea}
第一册实践卷的目标不是把 ELF、ABI、汇编和 benchmark 各讲一次，而是让它们在同一个工程里互相校验：字节是事实，解析是解释，测试固定合同，工具输出交叉验证，报告只写证据能支撑的结论。
\end{keyidea}
